% Sprint #1 — Panorama da Aviação Doméstica Brasileira
% Compilar no Overleaf com pdflatex
\documentclass[aspectratio=169]{beamer}

\usetheme{Madrid}
\setbeamertemplate{navigation symbols}{}
\setbeamertemplate{footline}[frame number]

\usepackage[utf8]{inputenc}
\usepackage[T1]{fontenc}
\usepackage[brazil]{babel}
\usepackage{booktabs}
\usepackage{tabularx}
\usepackage{tikz}
\usepackage{pgfplots}
\pgfplotsset{compat=1.18}
\usepackage{xcolor}
\usepackage{colortbl}
\usepackage{url}

% ── Cores ────────────────────────────────────────────────────────────────────
\definecolor{darkgray}{RGB}{80,80,80}
\definecolor{anacblue}{RGB}{0,82,147}
\definecolor{anacgold}{RGB}{200,150,0}
\definecolor{lightgray}{RGB}{240,240,240}
\definecolor{latamblue}{RGB}{55,138,221}
\definecolor{golgreen}{RGB}{99,153,34}
\definecolor{azulpurple}{RGB}{127,119,221}
\definecolor{passaredoamber}{RGB}{239,159,39}
\definecolor{othergray}{RGB}{136,135,128}

% ── Tema ─────────────────────────────────────────────────────────────────────
\setbeamercolor{palette primary}{bg=anacblue, fg=white}
\setbeamercolor{palette secondary}{bg=anacblue!80, fg=white}
\setbeamercolor{palette tertiary}{bg=anacblue!60, fg=white}
\setbeamercolor{structure}{fg=anacblue}
\setbeamercolor{block title}{bg=anacblue, fg=white}
\setbeamercolor{alerted text}{fg=anacgold}
\setbeamerfont{block title}{size=\small}

% ── Metadados ─────────────────────────────────────────────────────────────────
\title[Aviação Doméstica Brasileira]{Panorama da Aviação Doméstica Brasileira}
\author{Fernando Souza da Silveira}
\institute{PPGC -- UFRGS}
\date{2026}

\begin{document}

\begin{frame}
  \titlepage
\end{frame}

% ============================================================
\section{Dataset}
% ============================================================

\begin{frame}[t]{Fontes de Dados}
  Dados abertos da \textbf{ANAC} (\texttt{sistemas.anac.gov.br/dadosabertos}):
  \medskip
  \begin{columns}[T]
    \begin{column}{0.31\textwidth}
      \begin{block}{Produção Aeronáutica}
        \small
        \textbf{1,08 M} linhas $\cdot$ 38 atributos\\
        2000--2026 $\cdot$ 354 MB\\
        Empresa $\times$ rota $\times$ mês
      \end{block}
    \end{column}
    \begin{column}{0.31\textwidth}
      \begin{block}{Atrasos e Cancelamentos}
        \small
        \textbf{972 arquivos} CSV\\
        Anexos I, II e III\\
        27 anos $\times$ 12 meses
      \end{block}
    \end{column}
    \begin{column}{0.31\textwidth}
      \begin{block}{Aeródromos}
        \small
        \textbf{6.882} aeródromos\\
        Cadastro nacional\\
        Coordenadas, pistas, OACI
      \end{block}
    \end{column}
  \end{columns}
  \medskip
  \begin{alertblock}{Foco}
    Voos \textbf{domésticos}, 2016--2026.
    Comparativo: \textbf{SBPA/POA} (Salgado Filho) vs.\ média nacional.
  \end{alertblock}
\end{frame}

% ── Classificação ────────────────────────────────────────────────────────────
\begin{frame}[t]{Classificação dos Dados (Munzner)}
  \begin{columns}[T]
    \begin{column}{0.46\textwidth}
      \textbf{Tipo:} Tabela plana (\textit{flat table})\\[5pt]
      \begin{tabular}{@{}ll@{}}
        \textbf{Chaves}  & empresa, rota, ano, mês \\[3pt]
        \textbf{Valores} & métricas operacionais \\
      \end{tabular}
    \end{column}
    \begin{column}{0.50\textwidth}
      \textbf{Atributos:}
      \begin{itemize}\setlength{\itemsep}{3pt}
        \item \textbf{Nominal:} empresa, aeroporto, UF, país
        \item \textbf{Ord.\ sequencial:} ano
        \item \textbf{Ord.\ cíclico:} mês
        \item \textbf{Quantitativo:} passageiros, decolagens, ASK, RPK, carga\ldots
      \end{itemize}
    \end{column}
  \end{columns}
  \smallskip
  \footnotesize
  \begin{tabular}{@{}lll@{}}
    \toprule
    \textbf{Atributo} & \textbf{Tipo} & \textbf{Exemplo} \\
    \midrule
    \rowcolor{lightgray}
    \texttt{EMPRESA\_SIGLA}      & Nominal          & \texttt{GLO}, \texttt{AZU} \\
    \texttt{ANO / MES}           & Seq. / Cíclico   & \texttt{2024 / 7} \\
    \rowcolor{lightgray}
    \texttt{PASSAGEIROS\_PAGOS}  & Quantitativo     & \texttt{18\,432} \\
    \texttt{DECOLAGENS}          & Quantitativo     & \texttt{214} \\
    \rowcolor{lightgray}
    \texttt{ASK / RPK}           & Quantitativo     & assentos/pax $\times$ km \\
    \bottomrule
  \end{tabular}
\end{frame}

% ============================================================
\section{Persona}
% ============================================================

\begin{frame}[t]{Usuário Fictício}
  \begin{block}{Alberto Santos, 41 anos --- Gerente de Operações, Aeroporto Salgado Filho (SBPA)}
    \begin{itemize}\setlength{\itemsep}{7pt}
      \item Precisa apresentar à diretoria o \alert{desempenho operacional do SBPA}
            após a recuperação das enchentes de 2024
      \item Quer comparar o aeroporto com a média nacional e com outros hubs regionais
      \item Familiaridade com planilhas; \textbf{não} é especialista em visualização
      \item Espera um painel HTML interativo, sem instalação de software
    \end{itemize}
  \end{block}
\end{frame}

% ============================================================
\section{Tarefas}
% ============================================================

\begin{frame}[t]{Tarefas}
  \begin{columns}[T]
    \begin{column}{0.46\textwidth}
      \textbf{Alto nível}
      \begin{enumerate}\setlength{\itemsep}{5pt}
        \item \textbf{Descobrir} padrões do setor ao longo do tempo
        \item \textbf{Comparar} SBPA com outros aeroportos e média nacional
        \item \textbf{Identificar} recuperação pós-enchentes (2025--2026)
        \item \textbf{Localizar} o SBPA no contexto geográfico nacional
      \end{enumerate}
    \end{column}
    \begin{column}{0.50\textwidth}
      \textbf{Nível intermediário}
      \begin{itemize}\setlength{\itemsep}{3pt}
        \item Evolução mensal de passageiros e decolagens
        \item Calcular \textit{load factor} = RPK\,/\,ASK
        \item Top-5 rotas mais movimentadas do SBPA
        \item Percentual de atrasos e cancelamentos por mês/ano
        \item Market share por companhia no SBPA
        \item Impacto do COVID (2020) e enchentes (2024)
      \end{itemize}
    \end{column}
  \end{columns}
\end{frame}

% ============================================================
\section{Visualizações}
% ============================================================

% ── Viz 1 ────────────────────────────────────────────────────────────────────
\begin{frame}[t]{Viz 1 --- Mapa com Bolhas}
  \begin{columns}[c]
    \begin{column}{0.50\textwidth}
      \centering
      \begin{tikzpicture}[scale=0.72]
        % Contorno do Brasil (esquemático)
        \draw[anacblue!30, fill=anacblue!10, thick, rounded corners=8pt]
          (0.5,0) -- (4.5,0.3) -- (5.5,1.5) -- (5.2,4.5) --
          (4,5.5) -- (2,5.2) -- (0.3,3.5) -- (0,1.5) -- cycle;
        % Aeroportos: área ∝ passageiros
        \filldraw[anacblue!80, opacity=0.88] (3.5,2.0) circle (0.72cm);
        \filldraw[anacblue!80, opacity=0.88] (3.3,1.6) circle (0.52cm);
        \filldraw[anacblue!80, opacity=0.88] (3.0,3.5) circle (0.43cm);
        \filldraw[anacgold,    opacity=0.92] (2.2,0.6) circle (0.36cm);
        \filldraw[anacblue!70, opacity=0.85] (4.4,2.5) circle (0.28cm);
        \filldraw[anacblue!50, opacity=0.80] (1.5,3.0) circle (0.18cm);
        % Rótulos
        \node[font=\tiny, white] at (3.5,2.0) {GRU};
        \node[font=\tiny, white] at (3.3,1.6) {CGH};
        \node[font=\tiny, white] at (3.0,3.5) {BSB};
        \node[font=\tiny\bfseries, anacgold!10!black] at (2.2,0.6) {POA};
        \node[font=\tiny, white] at (4.4,2.5) {GIG};
        % Slider de ano (dentro do mapa, na parte inferior)
        \fill[white, opacity=0.85, rounded corners=2pt]
          (0.2,0.08) rectangle (5.3,0.44);
        \fill[anacblue!30, rounded corners=1pt]
          (0.4,0.18) rectangle (5.1,0.34);
        \filldraw[anacblue] (2.9,0.26) circle (0.10cm);
        \node[font=\tiny, darkgray] at (0.7,0.10) {\tiny 2016};
        \node[font=\tiny, darkgray] at (2.75,0.10) {\tiny Ano};
        \node[font=\tiny, darkgray] at (4.8,0.10) {\tiny 2026};
      \end{tikzpicture}
    \end{column}
    \begin{column}{0.47\textwidth}
      \begin{itemize}\setlength{\itemsep}{7pt}
        \item Área do círculo $\propto$ passageiros pagos
        \item \textcolor{anacgold}{\textbf{Dourado}} = SBPA em destaque
        \item Slider interativo para filtrar por ano
      \end{itemize}
      \medskip
      \textbf{Responde:} o SBPA é um hub regional ou figura entre os maiores aeroportos do país?
      \medskip\\
      \footnotesize\textbf{Dados:} \texttt{Dados\_Estatisticos.csv}\\
      + cadastro aeródromos (lat/lon)
    \end{column}
  \end{columns}
\end{frame}

% ── Viz 2 ────────────────────────────────────────────────────────────────────
\begin{frame}[t]{Viz 2 --- Série Temporal}
  \begin{columns}[c]
    \begin{column}{0.57\textwidth}
      \centering
      \begin{tikzpicture}
        \begin{axis}[
          width=7.8cm, height=4.2cm,
          scale only axis,
          xlabel={\footnotesize Ano},
          ylabel={\footnotesize Passageiros (M/mês)},
          xmin=2015.5, xmax=2026.5, ymin=0, ymax=12,
          grid=both,
          grid style={line width=0.3pt, draw=gray!25},
          tick label style={font=\tiny},
          label style={font=\footnotesize},
          xtick={2016,2018,2020,2022,2024,2026},
          ytick={0,3,6,9,12},
          legend style={font=\tiny, at={(0.99,0.05)}, anchor=south east,
                        draw=gray!40, fill=white, inner sep=3pt},
          axis line style={draw=gray!60},
        ]
          \addplot[anacblue, thick, smooth] coordinates {
            (2016,8.5)(2017,9.2)(2018,9.8)(2019,10.5)
            (2020,3.2)(2021,5.8)(2022,9.0)(2023,10.2)
            (2024,6.5)(2025,8.8)(2026,9.5)
          };
          \addplot[anacgold, thick, smooth, dashed] coordinates {
            (2016,0.8)(2017,0.9)(2018,1.0)(2019,1.1)
            (2020,0.3)(2021,0.6)(2022,0.9)(2023,1.1)
            (2024,0.2)(2025,0.7)(2026,0.9)
          };
          % Marcos de eventos
          \draw[red!55, dashed, thin] (axis cs:2020,0) -- (axis cs:2020,12);
          \node[font=\tiny, red!65, rotate=90, anchor=south]
            at (axis cs:2019.75,6.5) {COVID};
          \draw[orange!70, dashed, thin] (axis cs:2024,0) -- (axis cs:2024,12);
          \node[font=\tiny, orange!80, rotate=90, anchor=south]
            at (axis cs:2023.75,5.0) {Enchentes};
          \addlegendentry{Nacional}
          \addlegendentry{SBPA}
        \end{axis}
      \end{tikzpicture}
    \end{column}
    \begin{column}{0.40\textwidth}
      \begin{itemize}\setlength{\itemsep}{7pt}
        \item Nacional vs.\ SBPA mês a mês
        \item Marcos: COVID (2020) e enchentes RS (2024)
        \item Curva de recuperação visível
      \end{itemize}
      \medskip
      \textbf{Responde:} as enchentes de 2024 impactaram o SBPA mais do que a média nacional? O aeroporto já voltou ao nível anterior?
      \medskip\\
      \footnotesize\textbf{Dados:} \texttt{Dados\_Estatisticos.csv}
    \end{column}
  \end{columns}
\end{frame}

% ── Viz 3 ────────────────────────────────────────────────────────────────────
\begin{frame}[t]{Viz 3 --- Heatmap de Atrasos}
  \begin{columns}[c]
    \begin{column}{0.54\textwidth}
      \centering
      \begin{tikzpicture}[scale=0.82]
        % Rótulo do eixo Y
        \node[font=\tiny, rotate=90, anchor=south] at (-0.62, 1.50) {Ano};
        % Grade de células: linhas = anos (0=2022 bottom), colunas = meses
        \foreach \y/\ano in {0/2022, 1/2023, 2/2024, 3/2025, 4/2026} {
          \node[font=\tiny, anchor=east] at (-0.10, \y*0.60+0.30) {\ano};
          \foreach \x in {0,...,11} {
            \pgfmathsetmacro{\intens}{
              (\y==2 && \x>=3 && \x<=7) ? 92 :
              (\x==6 || \x==7)          ? 68 :
              (\x==11 || \x==0)         ? 58 :
              22 + \x*4
            }
            \fill[anacblue!\intens] (\x*0.50, \y*0.60) rectangle ++(0.47,0.56);
          }
        }
        % Rótulos dos meses
        \foreach \x/\m in
          {0/J,1/F,2/M,3/A,4/M,5/J,6/J,7/A,8/S,9/O,10/N,11/D}{
          \node[font=\tiny] at (\x*0.50+0.235, -0.20) {\m};
        }
        \node[font=\tiny] at (2.875, -0.42) {Mês};
        % Barra de cores abaixo (horizontal)
        \foreach \i in {0,...,8}{
          \pgfmathsetmacro{\v}{\i*11 + 5}
          \fill[anacblue!\v] (\i*0.635, -0.98) rectangle ++(0.60,0.28);
        }
        \node[font=\tiny, anchor=west]  at (0.00, -1.30) {0\% atrasos};
        \node[font=\tiny, anchor=east]  at (5.715,-1.30) {100\%};
      \end{tikzpicture}
    \end{column}
    \begin{column}{0.43\textwidth}
      \begin{itemize}\setlength{\itemsep}{7pt}
        \item Célula = mês $\times$ ano no SBPA
        \item Cor mais escura = maior \% atrasos $>$30 min
        \item 2024 (abr--ago) visivelmente atípico
      \end{itemize}
      \medskip
      \textbf{Responde:} os atrasos de 2024 foram consequência das enchentes ou o SBPA já mostrava piora antes do evento?
      \medskip\\
      \footnotesize\textbf{Dados:} \texttt{atrasos/*/anexo\_ii.csv}
    \end{column}
  \end{columns}
\end{frame}

% ── Viz 4 ────────────────────────────────────────────────────────────────────
\begin{frame}[t]{Viz 4 --- Market Share e Pontualidade}
  \begin{columns}[c]
    \begin{column}{0.58\textwidth}
      \centering
      \begin{tikzpicture}
        \begin{axis}[
          width=7.6cm, height=4.2cm,
          scale only axis,
          xlabel={\footnotesize\textit{Market share (\%)}},
          ylabel={\footnotesize\textit{Pontualidade (\%)}},
          xmin=-2, xmax=56,
          ymin=82, ymax=97,
          grid=none,
          tick label style={font=\tiny},
          label style={font=\footnotesize},
          xtick={0,10,20,30,40,50},
          ytick={83,86,89,92,95},
          axis line style={draw=gray!60},
          clip=false,
        ]
          % Quadrantes (fundos)
          \fill[passaredoamber, opacity=0.07]
            (axis cs:-2,89.2) rectangle (axis cs:19.6,97);
          \fill[latamblue, opacity=0.07]
            (axis cs:19.6,82) rectangle (axis cs:56,89.2);
          % Linhas de quadrante
          \draw[dashed, gray!45, thin]
            (axis cs:19.6,82) -- (axis cs:19.6,97);
          \draw[dashed, gray!45, thin]
            (axis cs:-2,89.2) -- (axis cs:56,89.2);
          % Rótulos de quadrante
          \node[font=\tiny\color{gray!70}, align=center] at (axis cs:8,96.0)
            {pontual\\+ nicho};
          \node[font=\tiny\color{gray!70}, align=center] at (axis cs:38,96.0)
            {pontual\\+ dominante};
          \node[font=\tiny\color{gray!70}, align=center] at (axis cs:8,83.2)
            {impont.\\+ nicho};
          \node[font=\tiny\color{gray!70}, align=center] at (axis cs:38,83.2)
            {impont.\\+ dominante};
          % Pontos
          \addplot[only marks, mark=*, mark size=4.5pt,
                   draw=latamblue, fill=latamblue]
            coordinates {(42,91)};
          \node[font=\tiny, right=2pt] at (axis cs:42,91) {LATAM};

          \addplot[only marks, mark=*, mark size=4.5pt,
                   draw=golgreen, fill=golgreen]
            coordinates {(31,87)};
          \node[font=\tiny, right=2pt] at (axis cs:31,87) {GOL};

          \addplot[only marks, mark=*, mark size=4.5pt,
                   draw=azulpurple, fill=azulpurple]
            coordinates {(19,89)};
          \node[font=\tiny, left=2pt] at (axis cs:19,89) {AZUL};

          \addplot[only marks, mark=*, mark size=5.5pt,
                   draw=passaredoamber, fill=passaredoamber]
            coordinates {(4,94)};
          \node[font=\tiny\bfseries\color{passaredoamber}, right=2pt]
            at (axis cs:4,94) {Passaredo};

          \addplot[only marks, mark=*, mark size=3.5pt,
                   draw=othergray, fill=othergray]
            coordinates {(2,85)};
          \node[font=\tiny\color{othergray}, right=2pt]
            at (axis cs:2,85) {Outras};
        \end{axis}
      \end{tikzpicture}
    \end{column}
    \begin{column}{0.39\textwidth}
      \begin{itemize}\setlength{\itemsep}{7pt}
        \item Eixo X = market share local
        \item Eixo Y = pontualidade média
        \item Quadrantes definidos pelas médias
        \item \textcolor{passaredoamber}{\textbf{Âmbar}} = destaque Passaredo
      \end{itemize}
      \medskip
      \textbf{Responde:} a empresa dominante no SBPA entrega a melhor experiência ao passageiro?
      \medskip\\
      \footnotesize\textbf{Dados:} \texttt{Dados\_Estatisticos.csv}\\
      \texttt{atrasos/*/anexo\_ii.csv}
    \end{column}
  \end{columns}
\end{frame}

% ============================================================
\section{Cronograma}
% ============================================================

\begin{frame}[t]{Cronograma}
  \footnotesize
  \begin{tabularx}{\textwidth}{@{} c l X c @{}}
    \toprule
    \textbf{Sprint} & \textbf{Período} & \textbf{Meta} & \textbf{Entrega} \\
    \midrule
    \rowcolor{lightgray}
    1 & 25/Mar -- 15/Abr
      & Especificação: dataset, persona, tarefas, 4~esboços
      & \textbf{15/Abr} \\
    2 & 16/Abr -- 06/Mai
      & ETL $\cdot$ Viz~1 (mapa com bolhas) $\cdot$ Viz~2 (série temporal)
      & \textbf{06/Mai} \\
    \rowcolor{lightgray}
    3 & 07/Mai -- 24/Jun
      & Viz~3 (heatmap) $\cdot$ Viz~4 (market share) $\cdot$ painel HTML
      & \textbf{24/Jun} \\
    \bottomrule
  \end{tabularx}
  \medskip\\
  \small\textbf{Stack:} Python 3.11 $\cdot$ Pandas $\cdot$ GeoPandas $\cdot$
  FastAPI $\cdot$ Next.js $\cdot$ D3.js
\end{frame}

% ============================================================
\section{Referências}
% ============================================================

\begin{frame}[t]{Referências}
  \footnotesize
  \begin{itemize}\setlength{\itemsep}{12pt}
    \item \textbf{Munzner, T.}
          \textit{Visualization Analysis and Design}.
          A K Peters/CRC Press, 2014.\\
          {\scriptsize\url{https://www.cs.ubc.ca/~tmm/vadbook/}}

    \item \textbf{ANAC} --- Agência Nacional de Aviação Civil.
          \textit{Dados Estatísticos} e \textit{Percentuais de Atrasos}.
          Acesso: Abr.\ 2026.\\
          {\scriptsize\url{https://sistemas.anac.gov.br/dadosabertos/Voos\%20e\%20opera\%C3\%A7\%C3\%B5es\%20a\%C3\%A9reas/}}
  \end{itemize}
\end{frame}

\begin{frame}[c]{}
  \centering
  {\LARGE\textbf{Obrigado!}}\\[10pt]
  \small\texttt{fernandosdba@gmail.com}
\end{frame}

\end{document}
